\section{Conclusion}
\label{sec:conclusion}
Modern software stacks integrate diverse input formats, from 3D models to graphics rendering and asynchronous network requests.
Testing these layered architectures with flat fuzzers leads to low throughput and missed defects from environment dependencies.
We presented \tool, an adaptive fuzzing architecture combining structure-aware mutation, a UCB1 bandit scheduler, and domain-specialized headless in-memory shims.
Across 59 target binaries in \fastgltf, \raylib, and \drogon, \tool executed 807{,}520{,}128 inputs and discovered 16 security vulnerabilities, comprising six upstream-resolved defects alongside ten additional medium- and high-severity findings.
Decoupling coverage accounting via linker wrapping and isolating subsystems into in-memory mocks enables high-throughput security testing across diverse software domains.
